\chapter{Espacios vectoriales}
\label{cap:espacios}

\section{Estructura vectorial}

Un \textbf{espacio vectorial sobre $\K$} es un conjunto no vacío $V$
con una suma de vectores y un producto de escalares por vectores.
La suma es asociativa y conmutativa, tiene un vector neutro $0$ y cada
vector $v$ tiene un opuesto $-v$. Además, para $u,v\in V$ y
$\alpha,\beta\in\K$,
\[
 \begin{gathered}
 \alpha(u+v)=\alpha u+\alpha v,\qquad
 (\alpha+\beta)v=\alpha v+\beta v,\\
 \alpha(\beta v)=(\alpha\beta)v,\qquad 1v=v.
 \end{gathered}
\]
El cero y los opuestos son únicos. Se deducen $0v=0$, $\alpha0=0$,
$(-1)v=-v$ y $\alpha v=0$ si y solo si $\alpha=0$ o $v=0$.
La suma permite cancelar un mismo vector en ambos miembros.

Son espacios vectoriales $\K^n$, $\M_{m\times n}(\K)$ y los polinomios
$\K[x]$, con sus operaciones habituales. También lo son
$\K_n[x]$, los polinomios de grado a lo sumo $n$ junto con el polinomio
nulo, y las funciones reales continuas con suma y producto puntual
por escalares. El cuerpo importa: $\C^n$ puede considerarse sobre
$\C$ o sobre $\R$, y su dimensión cambia.

\section{Combinaciones, independencia y generación}
\label{sec:generacion}

Una \textbf{combinación lineal} de $v_1,\ldots,v_m$ es
$\alpha_1v_1+\cdots+\alpha_mv_m$, con $\alpha_i\in\K$.
La familia es \textbf{linealmente independiente} si
\[
 \sum_{i=1}^m\alpha_iv_i=0\quad\Longrightarrow\quad
 \alpha_1=\cdots=\alpha_m=0.
\]
Es dependiente si hay una relación con algún coeficiente no nulo;
equivale a que alguno de sus vectores sea combinación de los demás.
Una familia que contiene $0$ es dependiente; un solo vector no nulo
es independiente. Dos vectores no nulos son dependientes si y solo si
son proporcionales.

Toda subfamilia de una familia independiente es independiente, y ampliar
una dependiente no elimina su dependencia. Añadir a una familia independiente
un vector conserva la independencia si y solo si el nuevo vector no es
combinación de los anteriores. Esto permite seleccionar sucesivamente
una subfamilia independiente que genere lo mismo que la original.

Se denota por $L(S)$ el conjunto de todas las combinaciones lineales
\emph{finitas} de vectores de $S$; para una lista se escribe
$L(v_1,\ldots,v_m)$. La familia es un \textbf{sistema generador} de $V$
si $L(v_1,\ldots,v_m)=V$. Se pueden añadir generadores o suprimir uno
que sea combinación de los demás sin cambiar el espacio generado.

Un espacio es de \textbf{dimensión finita} si admite un sistema generador
finito. En $\K[x]$ no existe tal sistema: cualquier familia finita
tiene grados acotados y no puede generar polinomios de todos los grados.
En lo sucesivo, los resultados numéricos sobre dimensión se refieren
a espacios de dimensión finita.

\section{Bases, dimensión y coordenadas}
\label{sec:bases}

Una \textbf{base} es una familia independiente que genera el espacio.
Equivale a un sistema generador mínimo o a una familia independiente
que no admite ampliación dentro del espacio. Para trabajar con coordenadas
se fija además un orden en la base.

Si una familia independiente tiene $r$ vectores y una generadora tiene $s$,
entonces $r\le s$. La idea es reemplazar sucesivamente generadores por
vectores independientes sin perder la generación. De aquí se obtiene
que todas las bases tienen el mismo número de vectores: la
\textbf{dimensión} del espacio, $\dim(V)$.

\begin{center}
\begin{tabular}{@{}lll@{}}
\toprule
Espacio & Base habitual & Dimensión \\
\midrule
$\K^n$ & $e_1,\ldots,e_n$ & $n$ \\
$\K_n[x]$ & $1,x,\ldots,x^n$ & $n+1$ \\
$\M_{m\times n}(\K)$ & Matrices $E_{ij}$ & $mn$ \\
$\{0\}$ & Familia vacía & $0$ \\
\bottomrule
\end{tabular}
\end{center}
En $E_{ij}$ solo la entrada $(i,j)$ vale $1$; las demás son cero.
Sobre $\R$, una base de $\C^n$ es $e_1,ie_1,\ldots,e_n,ie_n$,
por lo que su dimensión real es $2n$.

Si $\dim(V)=n$, toda familia independiente tiene a lo sumo $n$ vectores
y toda generadora, al menos $n$. Una familia de exactamente $n$ vectores
es base si y solo si es independiente, y si y solo si genera $V$.
De cualquier sistema generador se puede extraer una base; cualquier familia
independiente puede extenderse hasta una base, añadiendo vectores fuera
del espacio ya generado.

\subsection*{Coordenadas respecto de una base}

Para una base ordenada $B=(v_1,\ldots,v_n)$, cada vector tiene una
expresión única
\[
 v=x_1v_1+\cdots+x_nv_n.
\]
Los escalares son sus coordenadas; se escribe $v=(x_1,\ldots,x_n)_B$.
Para el cálculo matricial usamos la \emph{columna}
$\coord{v}{B}=(x_1,\ldots,x_n)^t$.
Las coordenadas respetan las operaciones:
\[
 \coord{\alpha u+\beta v}{B}
 =\alpha\coord{u}{B}+\beta\coord{v}{B}.
\]
La correspondencia entre $v$ y sus coordenadas es lineal y biyectiva.
Cambiar el orden de la base cambia el orden de las coordenadas.

\section{Rango de una familia y extracción de bases}
\label{sec:rango-vectores}

El \textbf{rango de una familia de vectores} es el máximo número de
vectores independientes que se pueden elegir de ella. Fijada una base $B$,
sea $C$ la matriz cuyas columnas son sus coordenadas. Entonces
\[
 \rg(v_1,\ldots,v_m)=\rg(C)=\dim L(v_1,\ldots,v_m).
\]
Colocar las coordenadas por filas da la matriz $C^t$, con el mismo rango.
Las operaciones elementales sobre la familia son reversibles y conservan
su rango y el subespacio generado.

Para obtener una base hay dos procedimientos útiles:
\begin{itemize}
 \item \textbf{Coordenadas por filas:} escalonar la matriz. Sus filas no nulas
       son coordenadas de una base del mismo subespacio. Los vectores así
       obtenidos pueden ser combinaciones de los originales; si se quieren
       identificar, se registran las mismas operaciones sobre sus nombres.
 \item \textbf{Coordenadas por columnas:} reducir por filas y localizar las
       columnas con pivote. Se eligen esas columnas de la matriz
       \emph{original}; corresponden a una base extraída de la familia dada.
\end{itemize}
Para calcular dependencias concretas se resuelve $C\Lambda=0$.
Los coeficientes de cada solución indican una relación entre los vectores.

\section{Cambio de base}
\label{sec:cambio-base}

Sean $B=(v_1,\ldots,v_n)$ y $B'=(w_1,\ldots,w_n)$ bases de $V$.
La \textbf{matriz de cambio de base de $B$ a $B'$} es
\[
 \M_{BB'}=\bigl(\coord{v_1}{B'}\ \cdots\ \coord{v_n}{B'}\bigr),
 \qquad \coord{x}{B'}=\M_{BB'}\coord{x}{B}.
\]
Sus columnas son los vectores de la base de partida expresados en la
base de llegada. Es invertible y
\[
 \M_{B'B}=\M_{BB'}^{-1},\qquad
 \M_{BD}=\M_{CD}\M_{BC}.
\]
El factor de la derecha realiza el primer cambio.
Si $S_B$ y $S_{B'}$ contienen las coordenadas de ambas bases por columnas
respecto de una misma base auxiliar, entonces
$\M_{BB'}=S_{B'}^{-1}S_B$.

\begin{ejemplo}
En $\R^2$, con $B=((1,1),(1,-1))$ y la base canónica $E$,
\[
 \M_{BE}=\begin{pmatrix}1&1\\1&-1\end{pmatrix},\qquad
 \M_{EB}=\frac12\begin{pmatrix}1&1\\1&-1\end{pmatrix}.
\]
El vector de coordenadas $(2,1)_B$ tiene coordenadas $(3,1)_E$.
La dirección del cambio determina cuál de las dos matrices se utiliza.
\end{ejemplo}

\section{Subespacios y sus ecuaciones}
\label{sec:ecuaciones-subespacios}

Un subconjunto \emph{no vacío} $U\subseteq V$ es un \textbf{subespacio}
si es cerrado para las combinaciones lineales:
\[
 u,v\in U,\quad\alpha,\beta\in\K
 \quad\Longrightarrow\quad\alpha u+\beta v\in U.
\]
Equivale a ser cerrado para suma y producto por escalares. Siempre contiene
el cero y tiene las operaciones heredadas de $V$. Los subespacios
$\{0\}$ y $V$ son el trivial y el total; los demás son propios.
Las dimensiones $1$, $2$ y $n-1$ corresponden a rectas, planos e hiperplanos
vectoriales, respectivamente. Un hiperplano se define por una ecuación
lineal homogénea no nula.

$L(S)$ es el menor subespacio que contiene $S$; si $S$ ya es un subespacio,
$L(S)=S$. Si $S\subseteq T$, entonces $L(S)\subseteq L(T)$.
Para comprobar igualdad de dos subespacios dados por generadores basta
mostrar que los generadores de cada uno pertenecen al otro.

\subsection*{Ecuaciones paramétricas e implícitas}

Fijada una base $B$ de $V$, sea $\dim(V)=n$ y $\dim(U)=d$.
Una base $u_1,\ldots,u_d$ de $U$ proporciona ecuaciones
\textbf{paramétricas}:
\[
 u=\sum_{i=1}^d\lambda_i u_i,\qquad
 \coord{u}{B}=C\Lambda,
 \quad C=(\coord{u_1}{B}\ \cdots\ \coord{u_d}{B}).
\]
Con una base, los $d$ parámetros son independientes y únicos.
Con generadores redundantes, la parametrización sigue siendo válida,
pero sus parámetros pueden no ser únicos.

Las ecuaciones \textbf{implícitas} describen $U$ mediante un sistema
homogéneo $D\coord{u}{B}=0$. Si se eliminan las ecuaciones redundantes,
hay exactamente $n-d$ ecuaciones independientes:
\[
 \dim(U)=n-\rg(D).
\]
Todo subespacio de $\K^n$ es el conjunto de soluciones de un sistema
lineal homogéneo, y todo sistema homogéneo define un subespacio.
En particular, las soluciones de $AX=0$ tienen dimensión $n-\rg(A)$.

\subsection*{Pasos para pasar de una descripción a otra}

\textbf{De generadores a implícitas:} se colocan sus coordenadas por columnas
en $C$ y se exige que $C\Lambda=X$ sea compatible, es decir,
$\rg(C)=\rg([\,C\mid X\,])$. Se escalona manteniendo $X$ simbólica:
las filas que tienen coeficientes nulos en $C$ dan las condiciones sobre $X$.
Después se eliminan condiciones redundantes. Otra posibilidad es partir
de una base de $U$, elegir un menor no nulo de orden $d$ de $C$ y anular
los menores que lo orlan con la columna $X$ y cada fila restante.
Para $U=\{0\}$, las ecuaciones son $X=0$; para $U=V$, no hay restricciones.

\textbf{De implícitas a paramétricas:} se resuelve $DX=0$ por Gauss,
asignando parámetros a las variables libres. Al dar a cada parámetro
el valor $1$ y a los demás $0$ se obtienen las coordenadas de una base.
Si no hay variables libres, la solución es solo $0$ y la base es vacía.

\section{Intersección, suma y suplementos}
\label{sec:suma-subespacios}

La intersección de subespacios es un subespacio: el mayor contenido en
todos ellos. Para obtener sus ecuaciones se reúnen las ecuaciones
implícitas de cada uno y se eliminan las redundantes.
La unión de subespacios no es, en general, un subespacio.

La \textbf{suma} es
\[
 U_1+\cdots+U_s=\{u_1+\cdots+u_s:u_i\in U_i\}
             =L(U_1\cup\cdots\cup U_s).
\]
Es el menor subespacio que los contiene a todos. Para hallar una base,
se concatenan sus generadores y se selecciona una familia independiente máxima.
Para calcular una intersección a partir de bases por columnas $C_U,C_W$,
se resuelve $C_U\alpha=C_W\beta$ y se forman los vectores comunes
$C_U\alpha$; después se selecciona una base de los obtenidos.

\subsection*{Suma directa}

La suma es \textbf{directa}, $U_1\oplus\cdots\oplus U_s$, si cada vector
de la suma admite una única descomposición con una componente en cada
subespacio. Son criterios equivalentes:
\[
 \begin{gathered}
 u_1+\cdots+u_s=0,\quad u_i\in U_i\ \Longrightarrow\ u_i=0\text{ para todo }i,\\
 U_i\cap\Bigl(\sum_{j\ne i}U_j\Bigr)=\{0\}\text{ para todo }i,\\
 \dim(U_1+\cdots+U_s)=\sum_{i=1}^s\dim(U_i).
 \end{gathered}
\]
En ese caso, concatenar bases de los sumandos da una base de la suma.
Para dos subespacios basta $U\cap W=\{0\}$; para tres o más,
las intersecciones por parejas no sustituyen al criterio con la suma
de los otros subespacios.

La \textbf{fórmula de Grassmann} relaciona las dimensiones:
\[
 \boxed{\dim(U+W)=\dim(U)+\dim(W)-\dim(U\cap W).}
\]
Su idea es extender una base de $U\cap W$ a bases de $U$ y de $W$:
al reunirlas, los vectores comunes se cuentan una sola vez.

$U$ y $W$ son \textbf{suplementarios en $V$} si $V=U\oplus W$.
Equivalentemente, $U+W=V$ y $U\cap W=\{0\}$.
Todo subespacio tiene un suplemento: se extiende una base suya a una
base de $V$ y se toman los vectores añadidos como base del suplemento.
Su dimensión es $\dim(V)-\dim(U)$ y, en general, no es único.
La igualdad de dimensiones por sí sola no garantiza que dos subespacios
sean suplementarios.

\section{Espacio cociente}
\label{sec:cociente}

Para un subespacio $U\subseteq V$, se define la relación de equivalencia
$v\sim_U w$ si $w-v\in U$. La clase de $v$ es
\[
 v+U=\{v+u:u\in U\}.
\]
El \textbf{espacio cociente} $V/U$ es el conjunto de estas clases.
Se tiene $v+U=w+U$ si y solo si $w-v\in U$, y $v+U=U$ si y solo si
$v\in U$. Sus elementos son clases de vectores, no representantes aislados.

Las operaciones son
\[
 (v+U)+(w+U)=(v+w)+U,\qquad \alpha(v+U)=\alpha v+U.
\]
Están bien definidas porque cambiar los representantes añade vectores
de $U$ y no cambia la clase resultante. El cero del cociente es la clase
$U=0+U$; el opuesto de $v+U$ es $-v+U$.

Si $(u_1,\ldots,u_d,v_{d+1},\ldots,v_n)$ es una base de $V$ que extiende
una base de $U$, entonces
\[
 (v_{d+1}+U,\ldots,v_n+U)\text{ es base de }V/U,
 \qquad \dim(V/U)=n-d.
\]
Las coordenadas de una clase se obtienen conservando las correspondientes
a los vectores añadidos: las componentes en $U$ desaparecen en el cociente.

\begin{ejemplo}
En $\R^2$, sea $U=L((1,0))$. La clase de $(a,b)$ es la recta horizontal
$\{(a+t,b):t\in\R\}$ y solo depende de $b$.
El cociente tiene dimensión $1$, con base $((0,1)+U)$; la coordenada
de $(a,b)+U$ en esa base es $b$. Su cero es toda la recta $U$.
\end{ejemplo}
